\section*{Chapter \ref{ch:fm:the-proprietary-market-making-firm} --- The Proprietary Market-Making Firm}

\begin{solution}{exo:fm:the-proprietary-market-making-firm:1}
Permanent minimum \$0.75 million; \omterm{def:fm:the-proprietary-market-making-firm:ownfunds}{fixed overheads requirement} $0.25\times95=\$23.75$ million; K-DTF $0.1\%\times10\,000=\$10$ million (no
net-position charge). The \omterm{def:fm:the-proprietary-market-making-firm:ownfunds}{fixed overheads requirement} binds: \$23.75 million.
\end{solution}

\begin{solution}{exo:fm:the-proprietary-market-making-firm:2}
K-DTF $0.01\%\times10\,000=\$1$ million. The requirement is still \$23.75 million, set by fixed overheads.
\end{solution}

\begin{solution}{exo:fm:the-proprietary-market-making-firm:3}
Allocations 46.5, 34.875, 23.25 and 11.625; draws \$37.2, 27.9, 18.6 and 9.3 million, \$93.0 million in all; \$23.25 million stays in the firm.
\end{solution}

\begin{solution}{exo:fm:the-proprietary-market-making-firm:4}
$g+d=35\%$ of the stock each year (\cref{prop:fm:the-proprietary-market-making-firm:parity}); spending grows at 15\%: $35\times1.15^9=\$123$
million in year 10.
\end{solution}

\begin{solution}{exo:fm:the-proprietary-market-making-firm:5}
Full-market revenue $10^{-4}\times20\,000\times250=\$500$ million. Profit before variable pay is zero when $35\times1.15^{t-1}=500-60=440$:
$t-1=\ln(440/35)/\ln1.15=18.1$, so year 20.
\end{solution}

\begin{solution}{exo:fm:the-proprietary-market-making-firm:6}
Above 20\% the firm fills its market in a few years and keeps retaining capital it cannot use, delaying payouts for no gain. Any fixed ratio either
retains too little while capital is short or too much once it is not; the fill-then-pay policy is optimal among capital-dependent policies
(\cref{prop:fm:the-proprietary-market-making-firm:retain}).
\end{solution}

\begin{solution}{exo:fm:the-proprietary-market-making-firm:7}
Year 56. Parity spending grows at 15\% a year and revenue at most at the market's 10\%, so spending catches up with revenue less fixed costs
eventually, however large the starting margin. A treadmill is sustainable only while the firm's market grows at least as fast as its competitors' spending.
\end{solution}

\begin{solution}{exo:fm:the-proprietary-market-making-firm:8}
The cut saves spending at once and loses capture only as the stock depreciates relative to competitors: in the model the lean firm earns more in
year 1 (\$124 million against \$116 million) and less from year 2, with capture at 0.44 of parity by year 10. Measure capture per unit traded and
relative speed, not one year's profit.
\end{solution}

\begin{solution}{pb:fm:the-proprietary-market-making-firm:1}
\begin{enumerate}
\item Regulatory capital, mostly equity and retained earnings; the requirement is the highest of the permanent minimum, a quarter of last year's
fixed overheads and the K-factors.
\item $0.25\times95=\$23.75$ million.
\item \$10 million on cash trades; \$1 million on derivatives notional.
\item The \omterm{def:fm:the-proprietary-market-making-firm:ownfunds}{fixed overheads requirement}, \$23.75 million; the prime brokers require \$200 million, 8.4 times it.
\item Net capital, net worth less illiquid assets and haircuts on positions, must stay above a floor set by the business and by indebtedness or
customers' debits.
\item A body corporate whose members share profit under an agreement and whose liability is limited; the capital a member contributed or left in,
repayable on departure.
\item \$10 billion a day, \$250 million, \$35 million, \$116.25 million.
\item Allocation \$46.5 million; draw \$37.2 million; \$9.3 million left in.
\item By the end of year 8, against year 6 without the departure.
\item The share of profit kept as capital.
\item \$582, 939, 706 and 296 million.
\item 20\%, \$939 million.
\item \$1\,210 million, \$271 million more: capital earns a high return until the market is full and nothing afterwards.
\item Years 6, 3 and 2.
\item Revenue depends on technology relative to competitors, so the firm must keep spending at their pace to stand still.
\item $s=(g+d)T$: 35\% of the stock.
\item 0.44 of parity.
\item Year 20 for the parity spender, year 18 for the lean spender.
\item Parity spending is 35\% of the technology stock, \$35 million in year 1 and \$123 million in year 10; retaining everything until capital fills
the market and then paying out all is worth \$1\,210 million against \$939 million for the best fixed ratio (20\%).
\item Keep profit until the firm's trading capital fills the market it can serve, then pay out; hold technology spending at the competitors' pace
and look for markets that grow at least as fast.
\end{enumerate}
\end{solution}

\begin{solution}{iq:fm:the-proprietary-market-making-firm:1}
Its \omterm{def:fm:the-proprietary-market-making-firm:ownfunds}{fixed overheads requirement} is a quarter of the previous year's fixed overheads: hiring and technology raise it whatever the firm trades.
\iqlookfor{the \omterm{def:fm:the-proprietary-market-making-firm:ownfunds}{fixed overheads requirement} and that it can bind before risk does.}
\end{solution}

\begin{solution}{iq:fm:the-proprietary-market-making-firm:2}
Usually the prime brokers' and clearing firms': they require capital posted against the positions, on their haircut and margin models, often
several times the regulatory figure.
\iqlookfor{posting requirements set trading capacity.}
\end{solution}

\begin{solution}{iq:fm:the-proprietary-market-making-firm:3}
Not all of it: while capital is below what the market can use, each retained dollar earns the firm's return on capital, far above the partners'
alternative. Keep enough to reach \$400 million, then pay out.
\iqlookfor{the marginal return on retained capital and the capacity cap.}
\end{solution}

\begin{solution}{iq:fm:the-proprietary-market-making-firm:4}
$g+d=35\%$ of the stock a year, and the amount grows at 10\%.
\iqlookfor{the steady-state condition $s=(g+d)T$.}
\end{solution}

\begin{solution}{iq:fm:the-proprietary-market-making-firm:5}
The departing capital is trading capital: repaying it at once would cut the firm's volume and could breach its requirements. Staggering turns the
withdrawal into a planned reduction.
\iqlookfor{capital as capacity, and the regulatory floor.}
\end{solution}

\begin{solution}{iq:fm:the-proprietary-market-making-firm:6}
Capture per unit traded, relative latency and fill quality over several quarters: the saving is immediate, the loss of capture arrives as the
technology falls behind competitors'.
\iqlookfor{the lag between spending and capture, and relative measures.}
\end{solution}
